%!TEX TS-program = xelatex

\documentclass[aspectratio=169]{beamer}

\newbool{russian}
\booltrue{russian}

% Подключаем стиль презентации HSE:
\usepackage{../last_year/HSE-theme/beamerthemeHSE}

% Работа с русским языком и шрифтами
\usepackage[english,russian]{babel}
\usepackage[no-math]{fontspec}
\setsansfont{HSESans-Regular}[
    Path = ../last_year/HSE_Sans/,
    BoldFont = HSESans-Bold.otf,
    ItalicFont = HSESans-Italic.otf
]
\setmonofont{Courier New}
\usepackage{mathspec}
\setmathsfont(Digits,Latin,Greek)[
  Path = ../last_year/HSE_Sans/,
  Numbers = {Lining,Proportional}
]{HSESans-Regular}
\setmathrm[Path = ../last_year/HSE_Sans/,Numbers={Lining,Proportional}]{HSESans-Regular}

\uselanguage{russian}
\languagepath{russian}
\deftranslation[to=russian]{Theorem}{Теорема}
\deftranslation[to=russian]{Definition}{Определение}
\deftranslation[to=russian]{Definitions}{Определения}
\deftranslation[to=russian]{Corollary}{Следствие}
\deftranslation[to=russian]{Fact}{Факт}
\deftranslation[to=russian]{Example}{Пример}
\deftranslation[to=russian]{Examples}{Примеры}

\usepackage{graphicx}
\usepackage{booktabs}
\usepackage{array}
\usepackage{colortbl}
\usepackage{listings}
\lstset{
  basicstyle=\tiny\ttfamily,
  backgroundcolor=\color{gray!12},
  breaklines=true,
  breakatwhitespace=false,
  frame=single,
  framesep=2pt,
  xleftmargin=2pt,
  xrightmargin=2pt,
  framerule=0.3pt,
  rulecolor=\color{gray!50}
}

%%%%%%%%%%%%%%%%%%%%%%%%%%%%%%%%%%%%%%%%%%%%%%%%%%%%%%%%%%%%
% Информация о выступлении

\title[Прогнозирование платёжных графов]{%
  \fontsize{11}{13}\selectfont Прогнозирование динамики графа платежей\\
  \fontsize{11}{13}\selectfont на основе открытых платёжных данных\\[0.2cm]
  \fontsize{9}{11}\selectfont Payment Graph Dynamics Forecasting\\
  \fontsize{9}{11}\selectfont Based on Open Payment Data
}

\subtitle{\fontsize{9}{11}\selectfont Командный программный проект}

\author[Кречетов, Романюков, Шеин]{%
  \fontsize{8}{10}\selectfont
  3 курс ОП «Прикладная математика и информатика»\\
  \textbf{Кречетов Роман Константинович} (\#БПМИ233),
  \textbf{Романюков Егор Сергеевич} (\#БПМИ233),
  \textbf{Шеин Илья Антонович} (\#БПМИ2310)\\[0.2cm]
  Руководитель проекта:\\
  \textbf{Филатов Дмитрий Станиславович},
  Исполнительный директор по исследованию данных, ПАО «Сбербанк»
}

\institute[НИУ ВШЭ]{%
  Национальный исследовательский университет\\
  «Высшая школа экономики»
}

\date{}
%%%%%%%%%%%%%%%%%%%%%%%%%%%%%%%%%%%%%%%%%%%%%%%%%%%%%%%%%%%%

\begin{document}

% Слайд 1: Титульный
\frame[plain]{\titlepage}

%%%%%%%%%%%%%%%%%%%%%%%%%%%%%%%%%%%%%%%%%%%%%%%%%%%%%%%%%%%%

% Слайд 2: Описание предметной области
\section{Описание предметной области}
\begin{frame}
  \frametitle{Описание предметной области}
  \begin{itemize}
    \item \textbf{Транзакционные графы} — естественное представление платёжных систем
      \begin{itemize}
        \item Узлы — участники системы: пользователи, организации, кошельки
        \item Рёбра — финансовые взаимодействия между ними с временными метками
      \end{itemize}
    \item \textbf{Bitcoin} — первая и наиболее популярная криптовалюта
      \begin{itemize}
        \item Все транзакции публично доступны с 2009 года
        \item Граф содержит сотни миллионов узлов и миллиарды рёбер
        \item Сильно разреженный, степени распределены по степенному закону
      \end{itemize}
    \item \textbf{Задача предсказания связей} (link prediction) — один из ключевых инструментов анализа подобных сетей
      \begin{itemize}
        \item Выявление мошеннических операций
        \item Прогнозирование транзакционной активности
        \item Анализ рыночных трендов и стратегическое планирование
      \end{itemize}
  \end{itemize}
\end{frame}

% Слайд 3: Актуальность работы
\section{Актуальность работы}
\begin{frame}[fragile]
  \frametitle{Актуальность работы}
  \begin{itemize}
    \item \textbf{Масштаб данных}: реальные финансовые сети превосходят стандартные бенчмарки на \textbf{1–4 порядка} по числу узлов и рёбер
    \item \textbf{Пробел в литературе}: большинство работ по temporal GNN тестируются на небольших социальных или биологических сетях; на крупных финансовых данных современные архитектуры проверялись значительно реже
    \item \textbf{Неизученный вопрос}: насколько эффективны простые структурные эвристики по сравнению с нейросетевыми моделями на реальных крупномасштабных финансовых данных?
    \item \textbf{Отсутствие воспроизводимого конвейера}: существующие библиотеки либо ориентированы на оффлайн-бенчмарки, либо не поддерживают форматы платёжных данных
  \end{itemize}
\end{frame}

% Слайд 4: Цель и задачи
\section{Цель и задачи}
\begin{frame}
  \frametitle{Цель и задачи}
  \textbf{Цель}
  \begin{itemize}
    \item Систематически сравнить методы предсказания временных связей в крупномасштабном транзакционном графе Bitcoin и разработать воспроизводимую платформу для экспериментов
  \end{itemize}
  \vspace{0.3cm}
  \textbf{Задачи}
  \begin{enumerate}
    \item Выбрать задачу и датасет; провести обзор существующих подходов и библиотек
    \item Разработать единый экспериментальный протокол (метрики, семплирование, признаки)
    \item Реализовать и сравнить $>$10 методов: структурные эвристики, ML-модели, Temporal GNN
    \item Провести аблационное исследование модели CN+MLP
    \item Разработать библиотеку \texttt{payment-graph-forecasting} с единым API
    \item Реализовать CUDA-модуль для ускорения вычислительно затратных примитивов
  \end{enumerate}
\end{frame}

% Слайд 5: Постановка задачи
\section{Постановка задачи}
\begin{frame}
  \frametitle{Постановка задачи: Temporal Link Prediction}
  \begin{columns}
    \column{0.55\textwidth}
    \textbf{Формальная постановка}
    \begin{itemize}
      \item Динамический граф $\mathcal{G} = (\mathcal{V}, \mathcal{E})$ — поток событий (stream graph)
      \item Каждое ребро $(s, d, t) \in \mathcal{E}$ — перевод от кошелька $s$ к $d$ в момент $t$
      \item Задача: по истории $\mathcal{G}_{<t}$ оценить вероятность $P(s, d, t \mid \mathcal{G}_{<t})$
    \end{itemize}
    \vspace{0.3cm}
    \textbf{Разбиение данных (temporal split)}
    \begin{itemize}
      \item Обучение: первые \textbf{70\%} рёбер
      \item Валидация: следующие \textbf{15\%}
      \item Тест: последние \textbf{15\%}
    \end{itemize}
    \column{0.42\textwidth}
    \textbf{Метрики качества}\\[0.2cm]
    \textbf{MRR} (Mean Reciprocal Rank):
    $$\text{MRR} = \frac{1}{|\mathcal{E}_\text{eval}|} \sum_{i=1}^{|\mathcal{E}_\text{eval}|} \frac{1}{r_i}$$
    \textbf{H@K} (Hits@K):
    $$\text{H@K} = \frac{1}{|\mathcal{E}_\text{eval}|} \sum_{i} \mathbf{1}[r_i \le K]$$
  \end{columns}
\end{frame}

% Слайд 6: Датасет
\section{Датасет}
\begin{frame}
  \frametitle{Датасет: ORBITAAL (Bitcoin, 2020)}
  \begin{columns}
    \column{0.5\textwidth}
    \scriptsize
    \textbf{ORBITAAL} — полный временной граф биткоин-транзакций (2009–2021)
    \begin{itemize}
      \item Используем 10\% рёбер периода \textbf{2020-06-01 — 2020-08-31}
      \item Соответствует $\approx$6.1M транзакций за 10 дней
    \end{itemize}
    \vspace{0.2cm}
    \begin{table}
    \small
    \begin{tabular}{lr}
      \toprule
      Характеристика & Значение \\
      \midrule
      Узлы $|\mathcal{V}|$ & 2\,638\,557 \\
      Рёбра $|\mathcal{E}|$ & 6\,143\,581 \\
      Уникальные пары & 5\,314\,320 \\
      Повторные взаимодействия & 829\,261 \\
      Плотность (directed) & $7.6 \times 10^{-7}$ \\
      Суммарный объём & \$45.4 млрд \\
      \bottomrule
    \end{tabular}
    \end{table}
    \column{0.47\textwidth}
    \footnotesize
    \textbf{Сравнение с классическими датасетами}
    \begin{table}
    \tiny
    \begin{tabular}{lrrl}
      \toprule
      Датасет & $|\mathcal{V}|$ & $|\mathcal{E}|$ & Период \\
      \midrule
      Contacts & 692 & 2.4M & 30 дн. \\
      LastFM & 1\,980 & 1.3M & 30 дн. \\
      Wikipedia & 9\,227 & 157K & 30 дн. \\
      Reddit & 10\,984 & 672K & 30 дн. \\
      AskUbuntu & 159K & 964K & 2613 дн. \\
      SuperUser & 194K & 1.4M & 2773 дн. \\
      \midrule
      \textbf{ORBITAAL (наш)} & \textbf{2.6M} & \textbf{6.1M} & \textbf{10 дн.} \\
      \bottomrule
    \end{tabular}
    \end{table}
    \vspace{0.2cm}
    {\small Наш датасет превосходит ближайший аналог (SuperUser) в \textbf{13.6 раза} по числу узлов}
  \end{columns}
\end{frame}

% Слайд 7: Анализ существующих библиотек
\section{Анализ существующих решений}
\begin{frame}
  \frametitle{Анализ существующих библиотек}
  \begin{itemize}
    \item \textbf{PyTorch Geometric Temporal} — ориентирован на снапшоты; ограниченная поддержка темпоральности «из коробки»
    \item \textbf{PyG (базовый)} — низкоуровневые примитивы \texttt{TemporalData}; пользователь собирает пайплайн самостоятельно
    \item \textbf{DGL} — мощный фреймворк общего назначения; без специализированной поддержки непрерывного времени
    \item \textbf{TGL (Amazon)} и \textbf{STM-Graph} — промышленные решения для сверхкрупных графов; плохо обобщаются за пределы своей ниши
    \item \textbf{DyGLib} — наиболее близкий ориентир: честный воспроизводимый бенчмарк для stream graph, охватывает несколько temporal методов
      \begin{itemize}
        \item Ограничения: CPU-only семплирование, привязка к фиксированному формату CSV/NPY, сложно расширять
      \end{itemize}
  \end{itemize}
  \vspace{0.2cm}
  \textbf{Вывод:} ни одна из существующих библиотек не является расширяемой платформой с поддержкой платёжных данных и GPU-ускорением семплирования
\end{frame}

% Слайд 8: Функциональные требования
\section{Функциональные требования}
\begin{frame}
  \frametitle{Функциональные требования}
  \begin{columns}
    \column{0.5\textwidth}
    \textbf{Требования к исследовательской части}
    \begin{itemize}
      \item Воспроизводимый протокол сравнения: единое разбиение, одинаковые метрики, одинаковое семплирование
      \item Поддержка $>$10 методов из четырёх классов
      \item Корректное построение признаков без утечки данных
      \item Три стратегии семплирования негативных рёбер (RND, HIST, IND)
    \end{itemize}
    \column{0.48\textwidth}
    \textbf{Требования к библиотеке}
    \begin{itemize}
      \item Единый библиотечный API (\texttt{payment\_graph\_forecasting.*})
      \item Декларативный запуск через YAML-конфигурации
      \item Поддержка формата stream graph (parquet)
      \item Взаимозаменяемые backend'ы: Python / C++ / CUDA
      \item Автоматическое тестирование (всё покрыто тестами)
      \item Документация и примеры запуска
    \end{itemize}
  \end{columns}
\end{frame}

% Слайд 9: Анализ методов — структурные и ML
\section{Анализ существующих методов}
\begin{frame}
  \frametitle{Методы предсказания связей: обзор}
  \begin{columns}
    \column{0.5\textwidth}
    \footnotesize
    \textbf{Структурные эвристики}
    \begin{itemize}
      \item \textbf{Common Neighbors (CN)}: $|\mathcal{N}(s) \cap \mathcal{N}(d)|$
      \item \textbf{Adamic–Adar (AA)}: взвешивает общих соседей обратно логарифму их степени
      \item \textbf{Jaccard}: нормирует CN на объединение окружений
      \item \textbf{Preferential Attachment (PA)}: $|\mathcal{N}(s)| \cdot |\mathcal{N}(d)|$
    \end{itemize}
    \vspace{0.2cm}
    \textbf{Классические ML-модели}
    \begin{itemize}
      \item \textbf{Logistic Regression} — линейная граница
      \item \textbf{Random Forest} — ансамбль деревьев
      \item \textbf{CatBoost} — градиентный бустинг
      \item \textbf{CN+MLP} — CN-признак + признаки узлов, двухслойный перцептрон
    \end{itemize}
    \column{0.47\textwidth}
    \footnotesize
    \textbf{Простые Temporal GNN-архитектуры}
    \begin{itemize}
      \item \textbf{EAGLE}: time-aware + structure-aware модули с T-PPR; линейная сложность
      \item \textbf{GraphMixer}: только MLP + mean-pooling, без механизма внимания
    \end{itemize}
    \vspace{0.2cm}
    \textbf{Трансформер-подобные Temporal GNN-архитектуры}
    \begin{itemize}
      \item \textbf{GLFormer}: адаптивная агрегация токенов с расширяющимся рецептивным полем;
      \item \textbf{DyGFormer}: трансформер над историческими последовательностями первых соседей;
    \end{itemize}
  \end{columns}
\end{frame}

% Слайд 10: Выбор методов
\begin{frame}
  \frametitle{Обоснование выбора методов}
  \begin{itemize}
    \item \textbf{Структурные эвристики} (CN, AA, Jaccard, PA): образуют сильный нижний порог; на реальных датасетах неожиданно конкурентоспособны с GNN \textit{[Poursafaei et al., 2022]}
    \item \textbf{Классические ML} (LogReg, RF, CatBoost, CN+MLP): проверяют линейную разделимость признаков и сочетание структурного сигнала с атрибутивной информацией
    \item \textbf{Нейросети}: выбраны 4 архитектуры, покрывающие широкий спектр:
      \begin{itemize}
        \item По вычислительной сложности: от почти линейных (EAGLE, GraphMixer) до квадратичных (DyGFormer)
        \item По способу учёта времени: адаптивные веса (EAGLE, GLFormer) vs. явный патчинг (DyGFormer, GraphMixer)
        \item По архитектуре: лёгкие агрегирующие (EAGLE, GraphMixer) vs. Transformer-подобные (GLFormer, DyGFormer)
      \end{itemize}
    \item Такой набор позволяет получить \textbf{более полную картину} без смещения в пользу одного класса методов
  \end{itemize}
\end{frame}

% Слайд 11: Признаки
\section{Особенности реализации}
\begin{frame}
  \frametitle{Признаковое описание узлов}
  \vspace{0.2cm}
  Каждый узел характеризуется \textbf{15 признаками}, вычисленными исключительно на обучающей выборке (без утечки данных):

  \vspace{0.2cm}
  \begin{columns}[T,onlytextwidth]
    \column{0.49\textwidth}
    \raggedright
    \textbf{Структура и взаимодействия}
    \begin{itemize}
      \item \textbf{Степенные}:\\
        $\log\,\text{in\_degree}$, $\log\,\text{out\_degree}$,\\
        \text{in\_out\_ratio}
      \item \textbf{Уникальные связанные адреса}:\\
        $\log\,\text{unique\_in\_cp}$,\\
        $\log\,\text{unique\_out\_cp}$
      \item \textbf{Энтропия связанных адресов}:\\
        out/in counterparty entropy
    \end{itemize}
    \column{0.49\textwidth}
    \raggedright
    \textbf{Объёмы и временная активность}
    \begin{itemize}
      \item \textbf{Транзакционные объёмы}:\\
        $\log\,\text{total\_btc\_in/out}$,\\
        $\log\,\text{avg\_btc\_in/out}$
      \item \textbf{Временные паттерны}:\\
        recency, activity\_span,\\
        $\log\,\text{event\_rate}$, burstiness
    \end{itemize}
  \end{columns}

  \vspace{0.1cm}
  {\footnotesize
    \begin{beamercolorbox}[wd=\textwidth,sep=0.10cm,rounded=true]{BlueOnWhite}
      \textbf{Пара $(s, d)$}: 15 признаков отправителя + 15 признаков получателя
      
      \textbf{Итог:} получается вектор размерности \textbf{30}
    \end{beamercolorbox}
  }
\end{frame}

% Слайд 12: Архитектура библиотеки
\section{Архитектура библиотеки}
\begin{frame}
  \frametitle{Архитектура библиотеки \texttt{payment-graph-forecasting}}
  \begin{columns}[T,onlytextwidth]
    \column{0.61\textwidth}
    {\scriptsize\bfseries Структура пакета}
    \vspace{0.03cm}
    \setlength{\fboxsep}{4pt}
    \noindent\colorbox{gray!8}{%
      \parbox[t][0.64\textheight][t]{\dimexpr\linewidth-2\fboxsep\relax}{%
        \raggedright
        \vspace*{2pt}
        {\fontsize{8.0}{7.5}\selectfont\ttfamily
        \renewcommand{\arraystretch}{1.05}
        \begin{tabular}{@{}l@{}}
        payment\_graph\_forecasting/ \\
        |-- config/ \\
        |-- data/ \\
        |-- models/ \\
        | |-- registry.py \\
        | `-- eagle.py, graphmixer.py, ... \\
        |-- experiments/ \\
        | `-- launcher.py, runners/ \\
        |-- training/ \\
        |-- evaluation/ \\
        |-- sampling/ \\
        |-- infra/ \\
        |-- analysis/ \\
        |-- cuda.py \\
        `-- graph\_metrics.py \\
        \end{tabular}}
      }%
    }
    \column{0.35\textwidth}
    \scriptsize
    \textbf{Как читать дерево}
    \begin{itemize}
      \setlength{\itemsep}{1pt}
      \item \textbf{Предметная область}: \texttt{data/}, \texttt{models/}, \texttt{sampling/}\\
      данные, адаптеры моделей и семплирование
      \item \textbf{Контур эксперимента}: \texttt{experiments/}, \texttt{training/}, \texttt{evaluation/}\\
      запуск, обучение и ranking-метрики
      \item \textbf{Системный слой}: \texttt{infra/}, \texttt{cuda.py}, \texttt{graph\_metrics.py}\\
      runtime-утилиты и ускоренные примитивы
    \end{itemize}
    \vspace{0.06cm}
    \textbf{Точки входа}
    \begin{itemize}
      \setlength{\itemsep}{1pt}
      \item \textbf{YAML}: \texttt{config} $\rightarrow$ \texttt{experiments.launcher}
      \item \textbf{Python API}: \texttt{models} + \texttt{training} + \texttt{evaluation}
    \end{itemize}
  \end{columns}
\end{frame}

% Слайд 13: YAML-пример
\begin{frame}[fragile]
  \frametitle{Использование библиотеки: YAML-конфигурация}
  \begin{columns}
    \column{0.51\textwidth}
\begin{lstlisting}
experiment:
  name: eagle_library_smoke
  model: eagle

data:
  source: stream_graph
  parquet_path: /tmp/stream_graph.parquet
  features_path: /tmp/features.parquet
  fraction: 0.10

sampling:
  strategy: tgb_mixed
  num_neighbors: 20
  n_random_neg: 50
  backend: auto

training:
  epochs: 5
  batch_size: 200
  lr: 0.001
\end{lstlisting}
    \column{0.46\textwidth}
    \footnotesize\textbf{Запуск одной командой:}
\begin{lstlisting}
python -m payment_graph_forecasting\
  .experiments.launcher \
  --config exps/examples/eagle.yaml
\end{lstlisting}
    \vspace{0.15cm}
    \footnotesize\textbf{Или через Python API:}
\begin{lstlisting}
from payment_graph_forecasting \
  import launch_experiment

spec = load_experiment_spec(
  "eagle.yaml")
result = launch_experiment(spec)
\end{lstlisting}
    \vspace{0.15cm}
    {\footnotesize GitHub: \texttt{paymentgraphforecasting/}\\
    \texttt{payment-graph-forecasting}}
  \end{columns}
\end{frame}

% Слайд 14: CUDA-модуль
\section{CUDA-модуль}
\begin{frame}
  \frametitle{CUDA-модуль: архитектура и примитивы}
  \small
  \begin{columns}[T]
    \column{0.56\textwidth}
    \textbf{Два реализованных примитива} (по 3 backend'а: Python / C++ / CUDA):
    \vspace{0.2cm}
    \begin{enumerate}
      \item \textbf{Поиск временных соседей} (\texttt{TemporalGraphSampler})\\
        Для узла $u$ и момента $t$: найти $K$ последних рёбер $(u, v, \tau)$ при $\tau < t$\\
        {\small CUDA: CSR-граф переносится в память GPU один раз; далее поиск соседей и извлечение признаков выполняются на GPU}
      \vspace{0.2cm}
      \item \textbf{Common Neighbors} (\texttt{CommonNeighbors})\\
        $\text{CN}(u, v) = |\mathcal{N}(u) \cap \mathcal{N}(v)|$\\
        {\small CUDA: bitset-представление множеств соседей и popcount для побитового пересечения}
    \end{enumerate}
    \column{0.40\textwidth}
    \textbf{Роль в библиотеке}
    \begin{itemize}
      \item \texttt{TemporalGraphSampler} используется как backend подготовки батча для GraphMixer, GLFormer и DyGFormer
      \item Стадия извлечения признаков собирает признаки соседей, рёбер и относительные временные интервалы
      \item \texttt{CommonNeighbors} используется как самостоятельная эвристика и как признак для ML-моделей
    \end{itemize}
  \end{columns}
\end{frame}

% Полная таблица бенчмарка TemporalGraphSampler
\begin{frame}
  \frametitle{Бенчмарк TemporalGraphSampler}
  \begin{table}
  \footnotesize
  \begin{tabular}{llrrrr}
    \toprule
    \textbf{Сценарий} & \textbf{Backend} & \textbf{Sample, мс} & \textbf{Feat, мс} & \textbf{Total, мс} & \textbf{Ускорение} \\
    \midrule
    graphmixer\_small & python & 2.74 & 4.82 & 7.55 & 1.0x \\
    graphmixer\_small & cpp    & 0.07 & 0.14 & 0.21 & 35.9x \\
    graphmixer\_small & cuda   & 0.08 & 0.07 & 0.15 & \textbf{49.9x} \\
    \midrule
    dygformer\_small  & python & 3.04 & 8.46 & 11.49 & 1.0x \\
    dygformer\_small  & cpp    & 0.40 & 3.63 & 4.04  & 2.8x \\
    dygformer\_small  & cuda   & 0.09 & 0.09 & 0.18  & \textbf{64.2x} \\
    \midrule
    dygformer\_medium & python & 14.10 & 48.30 & 62.33 & 1.0x \\
    dygformer\_medium & cpp    & 2.12  & 32.07 & 34.22 & 1.8x \\
    dygformer\_medium & cuda   & 0.12  & 0.22  & 0.34  & \textbf{183.7x} \\
    \midrule
    dygformer\_large  & python & 14.20 & 49.16 & 63.53 & 1.0x \\
    dygformer\_large  & cpp    & 2.14  & 33.08 & 35.26 & 1.8x \\
    dygformer\_large  & cuda   & 0.13  & 0.23  & 0.36  & \textbf{175.9x} \\
    \midrule
    real\_scale       & python & 2.79  & 5.89  & 8.64  & 1.0x \\
    real\_scale       & cpp    & 0.38  & 2.21  & 2.61  & 3.3x \\
    real\_scale       & cuda   & 0.09  & 0.09  & 0.18  & \textbf{48.5x} \\
    \bottomrule
  \end{tabular}
  \end{table}
  {\small Синтетические данные; ускорение — относительно python-backend}
\end{frame}

% Слайд 16: Common Neighbors benchmark
\begin{frame}
  \frametitle{Бенчмарк Common Neighbors}
  \begin{table}
  \footnotesize
  \begin{tabular}{llrrr}
    \toprule
    \textbf{Сценарий} & \textbf{Backend} & \textbf{Median, мс} & \textbf{p95, мс} & \textbf{Ускорение} \\
    \midrule
    sparse\_like\_bitcoin & python & 0.42 & 0.48 & 1.0x \\
    sparse\_like\_bitcoin & cpp    & 0.07 & 0.08 & \textbf{5.9x} \\
    sparse\_like\_bitcoin & cuda   & 1.18 & 1.24 & 0.4x \\
    \midrule
    breakeven            & python & 1.27 & 1.36 & 1.0x \\
    breakeven            & cpp    & 0.57 & 0.62 & 2.2x \\
    breakeven            & cuda   & 0.73 & 0.81 & 1.7x \\
    \midrule
    dense\_small         & python & 2.00 & 2.08 & 1.0x \\
    dense\_small         & cpp    & 1.02 & 1.06 & 1.9x \\
    dense\_small         & cuda   & 0.42 & 0.48 & 4.8x \\
    \midrule
    dense\_medium        & python & 8.51 & 8.88 & 1.0x \\
    dense\_medium        & cpp    & 5.05 & 5.21 & 1.7x \\
    dense\_medium        & cuda   & 0.71 & 0.75 & 12.0x \\
    \midrule
    dense\_large         & python & 35.79 & 37.14 & 1.0x \\
    dense\_large         & cpp    & 20.85 & 22.40 & 1.7x \\
    dense\_large         & cuda   & 1.66 & 1.71 & \textbf{21.5x} \\
    \bottomrule
  \end{tabular}
  \end{table}
  {\small \texttt{breakeven} — переходный режим, в котором CUDA уже начинает обгонять python, но ещё не становится лучшим backend'ом}
\end{frame}

% Слайд 17: CUDA — реальные данные
\begin{frame}
  \frametitle{CUDA-ускорение DyGFormer на реальных данных}
  \begin{columns}
    \column{0.5\textwidth}
    \small
    \textbf{Оптимизация}:
    \begin{itemize}
      \item Узкое место стандартного пайплайна DyGFormer — \textbf{подготовка батча}: темпоральное семплирование соседей, извлечение признаков и передача CPU$\to$GPU
      \item Решение: перенос этапа батч-подготовки на \texttt{TemporalGraphSampler} с backend \texttt{cuda}, встроенный в библиотечный API
      \item Активируется одной строкой конфига: \texttt{sampling.backend: cuda}
    \end{itemize}
    \column{0.47\textwidth}
    \small
    \textbf{Время обучения одной эпохи (GPU NVIDIA Tesla V100)}
    \begin{table}
    \small
    \begin{tabular}{lr}
      \toprule
      Конфигурация & Время, мин \\
      \midrule
      До оптимизации (CPU) & $\approx 55$ \\
      После (CUDA batch prep) & $\approx 27.5$ \\
      \midrule
      \textbf{Ускорение} & \textbf{$\approx$2x} \\
      \bottomrule
    \end{tabular}
    \end{table}
    \vspace{0.2cm}
    {\small Среднее по 11 измерениям: $1651.0$ с $\approx 27.5$ мин}\\
    \vspace{0.2cm}
    {\small Это \textbf{не синтетический бенчмарк} — ускорение проявляется в реальном обучающем пайплайне на реальных данных}
  \end{columns}
\end{frame}

% Слайд 16: Выбор средств реализации
\section{Выбор средств реализации}
\begin{frame}
  \frametitle{Выбор средств реализации}
  \begin{columns}
    \column{0.5\textwidth}
    \textbf{Основной стек}
    \begin{itemize}
      \item \textbf{Python 3.10+} — основной язык разработки
      \item \textbf{PyTorch / PyG} — реализация нейросетевых моделей
      \item \textbf{Pandas / Polars / PyArrow} — работа с данными в parquet
      \item \textbf{Pydantic} — типизированные конфигурации (\texttt{ExperimentSpec})
      \item \textbf{YAML} — декларативные конфигурации экспериментов
    \end{itemize}
    \column{0.47\textwidth}
    \textbf{Вычислительное ускорение}
    \begin{itemize}
      \item \textbf{NumPy / SciPy} — Python-backend для семплирования и CN
      \item \textbf{C++ extension (pybind11)} — cpp-backend: sorted merge, CSR-операции
      \item \textbf{CUDA / cuBLAS} — cuda-backend: битсеты, параллельный поиск соседей
    \end{itemize}
    \vspace{0.2cm}
    \textbf{Инфраструктура}
    \begin{itemize}
      \item \textbf{pytest} + всё покрыто тестами
      \item \textbf{GPU}: NVIDIA Tesla V100 для нейросетевых моделей
    \end{itemize}
  \end{columns}
\end{frame}

% Слайд 17: Результаты экспериментов
\section{Результаты экспериментов}
\begin{frame}
  \frametitle{Результаты сравнительного эксперимента}
  \begin{columns}[T]
    \column{0.5\textwidth}
    \centering
    \small
    \begin{tabular}{lcccc}
      \toprule
      \textbf{Модель} & \textbf{MRR} & \textbf{val MRR} & \textbf{H@1} & \textbf{H@10} \\
      \midrule
      \textbf{CN} & \textbf{0.64} & 0.70 & 0.55 & 0.81 \\
      AA & 0.63 & 0.69 & 0.54 & 0.79 \\
      CN+MLP & 0.62 & 0.69 & 0.53 & 0.79 \\
      Jaccard & 0.55 & 0.63 & 0.45 & 0.75 \\
      RF & 0.52 & 0.56 & 0.41 & 0.74 \\
      PA & 0.50 & 0.53 & 0.41 & 0.69 \\
      CatBoost & 0.49 & 0.54 & 0.38 & 0.73 \\
      LogReg & 0.35 & 0.36 & 0.26 & 0.52 \\
      \midrule
      EAGLE & 0.46 & 0.53 & 0.35 & 0.67 \\
      DyGFormer & 0.40 & 0.45 & 0.31 & 0.60 \\
      GraphMixer & 0.22 & 0.25 & 0.13 & 0.40 \\
      GLFormer & 0.09 & 0.07 & 0.11 & 0.20 \\
      \bottomrule
    \end{tabular}
    \column{0.47\textwidth}
    \textbf{Ключевые наблюдения}
    \begin{itemize}
      \item \textbf{CN достигает MRR = 0.64} — лучший результат среди всех методов
      \item Структурные эвристики \textbf{превосходят все GNN}: CN и AA опережают EAGLE на 0.18, DyGFormer на 0.24
      \item CN+MLP с одним CN-признаком близок к чистому CN — MLP восстанавливает структурный сигнал
      \item \textbf{Вывод}: в транзакционном графе Bitcoin структурная близость (наличие общих контрагентов) — доминирующий предиктор будущих связей
      \item Темпоральные паттерны, на которых специализируются GNN, менее информативны на этом датасете
    \end{itemize}
  \end{columns}
\end{frame}

% Слайд 18: Аблационное исследование CN+MLP
\begin{frame}
  \frametitle{Аблационное исследование CN+MLP}
  \begin{columns}[T]
    \column{0.55\textwidth}
    \centering
    \small
    \begin{tabular}{llcc}
      \toprule
      \textbf{ID} & \textbf{Признаки / условие} & \textbf{MRR} & \textbf{H@1} \\
      \midrule
      E6\_lr4 & cn\_uu, сниженный LR & \textbf{0.616} & \textbf{0.533} \\
      E5 & cn\_uu, cosine scheduler & 0.606 & 0.525 \\
      E1 & cn\_uu (baseline) & 0.603 & 0.522 \\
      E9\_big & cn\_uu, большая MLP & 0.597 & 0.518 \\
      E7\_e2 & cn\_uu + aa\_uu, короткое & 0.578 & 0.491 \\
      E2 & cn\_uu + aa\_uu + log1p\_cn & 0.572 & 0.487 \\
      E8\_all & 7 признаков, без dropout & 0.330 & 0.248 \\
      E4 & 7 признаков, BCE loss & 0.326 & 0.254 \\
      E3 & 7 признаков, BPR loss & 0.308 & 0.230 \\
      \bottomrule
    \end{tabular}
    \column{0.03\textwidth}
    \column{0.39\textwidth}
    \textbf{Выводы аблационного исследования}
    \begin{itemize}
      \item \textbf{cn\_uu — единственный значимый признак}: E1 (1 признак) vs. E2 (3 признака) — добавление ухудшает MRR $0.603 \to 0.572$
      \item \textbf{7 признаков резко снижают качество}: E3–E8 дают MRR $\approx 0.31$–$0.33$ — примерно вдвое хуже E1
      \item \textbf{Увеличение MLP не помогает}: E9\_big $\approx$ E1; узкое место — входной скалярный признак, а не ёмкость модели
      \item \textbf{Тюнинг LR важнее архитектуры}: лучший результат E6\_lr4 достигнут без изменения признаков
    \end{itemize}
  \end{columns}
\end{frame}

% Слайд 19: Основные результаты
\section{Результаты и выводы}
\begin{frame}
  \frametitle{Основные результаты}
  \begin{enumerate}
    \item \textbf{Сравнительный анализ}: протестировано $>$10 методов четырёх классов на едином разбиении реального Bitcoin-датасета по метрикам MRR, H@1, H@10
    \vspace{0.2cm}
    \item \textbf{Превосходство структурных методов}: CN достигает MRR = 0.64 и превосходит все GNN-архитектуры — наличие общих контрагентов является доминирующим предиктором на этом датасете
    \vspace{0.2cm}
    \item \textbf{Аблационное исследование CN+MLP}: единственный значимый признак — cn\_uu; расширение признакового пространства стабильно ухудшает результат; наибольший прирост — от тюнинга LR
    \vspace{0.2cm}
    \item \textbf{Библиотека payment-graph-forecasting}: расширяемая платформа с единым API, YAML-конфигурациями, 7 моделями, 47 тестами и документацией
    \vspace{0.2cm}
    \item \textbf{CUDA-ускорение}: до 183x ускорения поиска временных соседей на синтетических бенчмарках; 2x ускорение обучения DyGFormer на реальных данных ($55 \to 27.5$ мин/эпоха)
  \end{enumerate}
\end{frame}

% Слайд 20: Направления дальнейшей работы
\section{Направления дальнейшей работы}
\begin{frame}
  \frametitle{Направления дальнейшей работы}
  \begin{columns}
    \column{0.5\textwidth}
    \textbf{Исследовательские направления}
    \begin{itemize}
      \item Проверить гипотезу на других финансовых датасетах (Ethereum, банковские данные)
      \item Разработать гибридные модели, объединяющие структурный сигнал с темпоральными паттернами
      \item Исследование более сложных стратегий семплирования (HIST, IND негативы)
    \end{itemize}
    \column{0.47\textwidth}
    \textbf{Инженерные направления}
    \begin{itemize}
      \item Расширение библиотеки: добавление новых моделей через адаптерный слой без изменения API
      \item Поддержка распределённого обучения (multi-GPU)
      \item Расширение CUDA-модуля: параллельный Common Neighbors для более плотных графов
      \item Интеграция с облачными хранилищами данных
      \item HPO (hyperparameter optimization) entrypoints в \texttt{experiments} модуле
    \end{itemize}
  \end{columns}
\end{frame}

% Слайд 21: Доступ к коду
\begin{frame}
  \frametitle{Доступ к исходному коду}
  \begin{itemize}
    \item \textbf{Проект доступен как open source на GitHub:}
      \begin{itemize}
        \item \url{https://github.com/paymentgraphforecasting/payment-graph-forecasting}
      \end{itemize}
    \item \textbf{В репозитории доступны:}
      \begin{itemize}
        \item Канонический библиотечный API (\texttt{payment\_graph\_forecasting/})
        \item 7 reference YAML-конфигураций в \texttt{exps/examples/}
        \item Документация в \texttt{docs/}: CUDA-модуль, evaluation protocol, сценарии работы
        \item 47 автоматических тестов в \texttt{tests/}
        \item Инструкции по установке и сборке CUDA-расширений
      \end{itemize}
  \end{itemize}
\end{frame}

% Слайд 22: Список источников
\section{Список источников}
\begin{frame}[allowframebreaks]
  \frametitle{Список использованных источников}
  \small
  \begin{thebibliography}{99}
    \bibitem{orbitaal} \textit{Béres F. et al.} ORBITAAL: A Comprehensive Bitcoin Dataset for Temporal Graph Analysis. — Zenodo, 2022. (CC BY 4.0)
    \bibitem{dyglib} \textit{Yu L. et al.} Towards Better Dynamic Graph Learning: New Architecture and Unified Library. — NeurIPS, 2023.
    \bibitem{tgb} \textit{Poursafaei F. et al.} Towards Better Evaluation for Dynamic Link Prediction. — NeurIPS, 2022.
    \bibitem{cn} \textit{Liben-Nowell D., Kleinberg J.} The Link Prediction Problem for Social Networks. — JACM, 2007.
    \bibitem{eagle} \textit{Li Z. et al.} EAGLE: Efficient and Effective GrAph modEL for temporal Link prEdiction. — 2024.
    \bibitem{graphmixer} \textit{Cong W. et al.} Do We Really Need Complicated Model Architectures for Temporal Networks? — ICLR, 2023.
    \bibitem{glformer} \textit{Kong L. et al.} GLFormer: Temporal Link Prediction with Global-Lens Transformer. — 2024.
    \bibitem{dygformer} \textit{Yu L. et al.} DyGFormer: A Graph Transformer for Dynamic Graphs. — NeurIPS, 2023.
    \bibitem{catboost} \textit{Prokhorenkova L. et al.} CatBoost: unbiased boosting with categorical features. — NeurIPS, 2018.
    \bibitem{ctdne} \textit{Nguyen G. H. et al.} Continuous-Time Dynamic Network Embeddings. — WWW Companion, 2018.
  \end{thebibliography}
\end{frame}

\end{document}
